\subsubsection*{命题 3.2.}
设 $X$ 是域 $k$ 上维数 $n\geq1$ 的光滑完全交. 若
\begin{equation}\tag{3.2.1}
H^n(\overline X,\mathbb Q_\ell)=N^1H^n(\overline X,\mathbb Q_\ell),
\end{equation}
则 $X$ 特殊; 亦即,
\begin{equation}\tag{3.2.2}
\mathfrak F_p\quad\text{在下述空间上幂零:}\quad H^n(X,\mathcal O_X).
\end{equation}

\medskip
\noindent\emph{证明.}
特征零的情形是平凡的.\EditorialFootnote{原书第381页写作“$p=0$ 的情形”, 此处的 $p$ 指域的特征. 若沿用 1.0 的“特征指数”约定, 特征零应写作 $p=1$. 下文均讨论正特征.}
故可设 $k$ 在 $\mathbb F_p$ 上有限生成, 且存在 $k$ 的有限生成
$\mathbb F_p$ 子代数 $B$, 使 $X/k$ 是光滑相对完全交 $\widetilde X$ 的泛纤维:
\begin{equation}\tag{3.2.3}
\begin{tikzcd}[column sep=4.5em,row sep=3em]
\widetilde X \arrow[d,"\pi"'] &
X=\widetilde X_\eta \arrow[l] \arrow[d] \\
\operatorname{Spec}(B)=S & \eta=\operatorname{Sp}(k)\ . \arrow[l]
\end{tikzcd}
\end{equation}
$B$ 模 $R^n\pi_*(\mathcal O_{\widetilde X})$ 局部自由, 其形成与任意基变换可交换.
因此点集
\begin{equation}\tag{3.2.4}
\{\,s\in S\mid X_s\text{ 特殊}\,\}
\end{equation}
在 $S$ 中闭.

由 (XX~2.3), 缩小 $S$ 后可设
\begin{equation}\tag{3.2.5}
\forall s\in S,\qquad
H^n(X_{\bar s},\mathbb Q_\ell)=N^1H^n(X_{\bar s},\mathbb Q_\ell).
\end{equation}
由推论~3.1, 对 $S$ 的每个闭点 $s$, $X_s$ 都特殊.
因此闭子集 (3.2.4) 等于整个 $S$, 从而特别地 $X=\widetilde X_\eta$ 特殊.

\medskip
\noindent\emph{底本注.}
图 (3.2.3) 中, 照录法文层保留了印本记号 $X=\widetilde X/_{\eta}$.
校订法文层和英文层使用标准的泛纤维记号 $X=\widetilde X_\eta$, 与证明末句一致.
